\documentclass[11pt,a4paper]{article}
\usepackage[utf8]{inputenc}
\usepackage[T1]{fontenc}
\usepackage[french]{babel}
\usepackage[a4paper,margin=2.2cm]{geometry}
\usepackage{lmodern}
\usepackage{xcolor}
\usepackage{hyperref}
\usepackage{enumitem}
\usepackage{tabularx}
\usepackage{booktabs}
\usepackage{array}
\usepackage{microtype}
\hypersetup{colorlinks=true,linkcolor=black,urlcolor=blue,citecolor=black}
\setlength{\parindent}{0pt}
\setlength{\parskip}{0.45em}
\renewcommand{\arraystretch}{1.15}
\definecolor{titleblue}{RGB}{35,74,108}

\begin{document}

{\Large\bfseries\color{titleblue} Informations du jour --- 1\textsuperscript{er} mai 2026}\par
{\small Synthèse vérifiée à partir de sources ouvertes consultées le 1\textsuperscript{er} mai 2026.}\vspace{0.6em}

\section*{Synthèse rapide}
\begin{itemize}[leftmargin=1.2em,itemsep=0.25em]
    \item \textbf{France -- vie quotidienne :} en mai 2026, le Gouvernement annonce l'extension progressive du repas à 1 euro à tous les étudiants à partir du 4 mai, une hausse moyenne du prix repère de vente du gaz de 15,4\,\%, et l'entrée en vigueur d'un code des douanes réorganisé au 1\textsuperscript{er} mai.
    \item \textbf{France -- économie :} l'économie française a stagné au premier trimestre 2026, en dessous de la prévision de croissance de 0,2\,\% mentionnée par Reuters à partir des estimations préliminaires.
    \item \textbf{Union européenne :} l'accord UE--Mercosur entre en application provisoire le 1\textsuperscript{er} mai 2026, malgré les réserves politiques, agricoles et environnementales, notamment en France.
    \item \textbf{Science et environnement :} la NASA signale comme image du jour des glissements de terrain en Papouasie-Nouvelle-Guinée après les pluies associées au cyclone tropical Maila.
\end{itemize}

\section*{France : mesures entrant en vigueur en mai}
\begin{tabularx}{\textwidth}{>{\bfseries}p{0.25\textwidth}X}
\toprule
Thème & Information essentielle \\
\midrule
Étudiants & À partir du 4 mai 2026, le repas à 1 euro est progressivement étendu à tous les étudiants dans les restaurants universitaires, notamment les titulaires d'une carte d'étudiant, d'une carte d'étudiant des métiers et les doctorants. \\
Gaz & À partir du 1\textsuperscript{er} mai 2026, le prix repère de vente du gaz augmente en moyenne de 15,4\,\%, soit 6,19 euros par mois, pour les offres indexées sur ce prix repère. \\
Douanes & Le code des douanes est réorganisé au 1\textsuperscript{er} mai 2026 : les règles ne changent pas sur le fond, mais la présentation est simplifiée pour les entreprises et les professionnels. \\
\bottomrule
\end{tabularx}

\textit{Lecture opérationnelle :} l'effet immédiat le plus concret pour les ménages est la hausse du prix repère du gaz ; pour l'enseignement supérieur, l'extension du repas à 1 euro est une mesure sociale importante susceptible d'avoir un impact direct sur les étudiants et doctorants.

\section*{Économie : stagnation française au premier trimestre 2026}
Selon Reuters, l'économie française est restée stable au premier trimestre 2026. Cette stagnation marque un ralentissement par rapport au trimestre précédent, où la croissance avait été de 0,2\,\%. La performance est inférieure aux attentes mentionnées par Reuters, qui faisaient état d'une prévision de 0,2\,\% pour le trimestre. L'explication donnée dans l'article est un effet de compensation : les contributions positives des stocks ont été contrebalancées par une contribution défavorable du commerce extérieur.

\textit{Point d'attention :} pour l'analyse macroéconomique, cette stagnation doit être lue avec prudence, car il s'agit d'une estimation préliminaire. Elle signale toutefois une fragilité de la dynamique française dans un contexte européen marqué par des tensions énergétiques et commerciales.

\section*{Europe et commerce : application provisoire de l'accord UE--Mercosur}
Reuters indique que l'Union européenne et le Mercosur commencent l'application provisoire de leur accord commercial à partir du 1\textsuperscript{er} mai 2026. Le texte reste politiquement sensible : ses soutiens, notamment l'Allemagne et l'Espagne, y voient un instrument de diversification commerciale, tandis que la France et d'autres critiques soulignent les risques pour les agriculteurs européens et pour l'environnement.

\textit{Enjeu stratégique :} l'accord vise à renforcer la capacité commerciale européenne dans un contexte de tensions avec les États-Unis et de concurrence chinoise accrue. Les gains macroéconomiques annoncés restent cependant modestes et l'effet complet ne peut être attendu qu'à long terme.

\section*{Science, climat et observation de la Terre}
La page Earth Science de la NASA met en avant, le 1\textsuperscript{er} mai 2026, un épisode de glissements de terrain en Papouasie-Nouvelle-Guinée provoqués par les fortes pluies associées au cyclone tropical Maila. La NASA rappelle aussi que sa division Earth Science exploite plus de vingt satellites et finance des programmes de recherche dédiés à l'observation de l'océan, de l'atmosphère, des terres, des glaces et des interactions entre ces compartiments.

\textit{Intérêt scientifique :} cet exemple illustre la pertinence de l'observation satellitaire pour relier aléas hydrométéorologiques, instabilités de terrain et gestion des risques naturels.

\section*{Sources principales}
\begin{enumerate}[leftmargin=1.5em,itemsep=0.2em]
    \item \href{https://www.info.gouv.fr/actualite/ce-qui-change-en-mai-2026}{Gouvernement français, \textit{Ce qui change en mai 2026}}, publié le 22/04/2026, modifié le 30/04/2026.
    \item \href{https://www.reuters.com/markets/us/french-economy-flat-first-quarter-below-forecasts-preliminary-gdp-2026-04-30/}{Reuters, \textit{French economy flat in the first quarter, below forecasts - preliminary GDP}}, 30/04/2026.
    \item \href{https://www.reuters.com/world/china/eu-kickstarts-mercosur-pact-counter-us-trade-hit-2026-04-30/}{Reuters, \textit{EU kickstarts Mercosur pact to counter US trade hit}}, 30/04/2026.
    \item \href{https://science.nasa.gov/earth/}{NASA Science, \textit{Earth}}, actualités et image du jour du 1\textsuperscript{er} mai 2026.
\end{enumerate}

\end{document}
